\section{复现命令}
\label{app:repro}

均在 \texttt{embdb/x} 下、WSL 环境：

\begin{lstlisting}[language=bash,caption={回归与关键探针}]
make -j2 -s build/g18 && ./build/g18
make -j2 -s test          # 全量较久
# layout 峰值（需先编 probe）
g++ -std=c++17 -O2 -pthread -Iinclude \
  bench/probe/layout.cpp -o build/layout
# det 口径示例：budget=0 units=8
./build/layout 200000 500 32 1 1 0 0 0 0 8
# 护洞 / 碎片：见 backup/r18/run_layout_*.sh
bash backup/peers-r17/run_peers_r17.sh      # R17 基线
bash backup/peers-soft/run_peers_soft.sh    # Soft 产品形态
# 或：bash backup/peers-soft/run_all.sh     # make test + Soft peers
\end{lstlisting}

\begin{lstlisting}[language=bash,caption={本文档编译}]
cd docs/latex
make          # xelatex ×2
# 产物：build/main.pdf
\end{lstlisting}

\subsection*{参考文献（仓库内）}
\begin{itemize}[leftmargin=1.6em]
  \item \texttt{docs/README.md} — Markdown 分类索引
  \item \texttt{docs/10-design/DESIGN3.md} — 设计与历史坑清单 §12
  \item \texttt{docs/00-entry/EVAL17.md} — 冻结全量评估
  \item \texttt{docs/20-rounds/ROUND15.md}--\texttt{docs/20-rounds/ROUND18.md} — 各轮执行与数字
  \item \texttt{docs/30-discuss/DISCUSS15\_FSYNC.md} — fsync 出锁讨论稿
  \item \texttt{docs/10-design/P2\_DESIGN.md} — 格式 v2 设计
  \item \texttt{docs/00-entry/BATTLE.md} — 升格与分叉对照
  \item \texttt{backup/README.md} — 归档 / peers 日志分类
\end{itemize}
